\chapter{2008年天津卷（理）}

\section{单选题}

\begin{problem}
\(\i\) 是虚数单位，\( \frac{\i^{3}(\i+1)}{\i-1}= \)（\(\quad\)）

\choices
    {\(-1\)}
    {\(1\)}
    {\(-\i\)}
    {\(\i\)}
\end{problem}

\begin{answer}
A
\end{answer}

\begin{solution}
因为 \(\i^3=-\i\)，所以原式为 \(\frac{-\i(\i+1)}{\i-1}=\frac{1-\i}{\i-1}=-1\).
故选 A.
\end{solution}

\begin{problem}
设变量 \(x\)，\(y\) 满足约束条件 \(\begin{cases}
x - y \geqslant 0,\\
x + y \leqslant 1,\\
x + 2y \geqslant 1,
\end{cases}\) 则目标函数 \(z = 5x + y\) 的最大值为（\(\quad\)）

\choices
    {\(2\)}
    {\(3\)}
    {\(4\)}
    {\(5\)}
\end{problem}

\begin{answer}
D
\end{answer}

\begin{solution}
由约束条件得可行域的边界分别为 \(y=x\)、\(y=1-x\) 和 \(y=\frac{1-x}{2}\). 三条边界直线两两相交于 \(\left(\frac12,\frac12\right)\)、\(\left(\frac13,\frac13\right)\) 和 \((1,0)\). 这三个交点均满足约束条件，构成可行域的三个顶点；在线性目标函数下，只需比较顶点处的值. 目标函数在这三个顶点处的值分别为 \(3\)、\(2\)、\(5\). 因此 \(z\) 的最大值为 \(5\)，故选 D.
\end{solution}

\begin{problem}
设函数 \(f(x) = \sin \left(2x - \frac{\pi}{2}\right)\)，\(x \in \R\)，则 \( f(x) \) 是（\(\quad\)）

\choices
    {最小正周期为 \(\pi\) 的奇函数}
    {最小正周期为 \(\pi\) 的偶函数}
    {最小正周期为 \(\frac{\pi}{2}\) 的奇函数}
    {最小正周期为 \(\frac{\pi}{2}\) 的偶函数}
\end{problem}

\begin{answer}
B
\end{answer}

\begin{solution}
由诱导公式，\(f(x)=\sin\left(2x-\frac{\pi}{2}\right)=-\cos 2x\).
函数 \(\cos 2x\) 的最小正周期为 \(\pi\)，且是偶函数，因此 \(f(x)\) 的最小正周期为 \(\pi\)，也是偶函数，故选 B.
\end{solution}

\begin{problem}
设 \(a\)，\(b\) 是两条直线，\(\alpha\)，\(\beta\) 是两个平面，则 \(a \perp b\) 的一个充分条件是（\(\quad\)）

\choices
    {\(a \perp \alpha\)，\(b \parallel \beta\)，\(\alpha \perp \beta\)}
    {\(a \perp \alpha\)，\(b \perp \beta\)，\(\alpha \parallel \beta\)}
    {\(a \subset \alpha\)，\(b \perp \beta\)，\(\alpha \parallel \beta\)}
    {\(a \subset \alpha\)，\(b \parallel \beta\)，\(\alpha \perp \beta\)}
\end{problem}

\begin{answer}
C
\end{answer}

\begin{solution}
若 \(a\subset\alpha\)，\(b\perp\beta\)，且 \(\alpha\parallel\beta\)，则垂直于平面 \(\beta\) 的直线 \(b\) 也垂直于平面 \(\alpha\). 因为直线 \(a\) 在平面 \(\alpha\) 内，所以 \(a\perp b\). 因此选项 C 给出了 \(a\perp b\) 的充分条件，故选 C.
\end{solution}

\begin{problem}
设椭圆 \(\frac{x^2}{m^2} + \frac{y^2}{m^2 - 1} = 1\)（\(m > 1\)）上一点 \(P\) 到其左焦点的距离为 \(3\)，到右焦点的距离为 \(1\)，则 \(P\) 点到右准线的距离为（\(\quad\)）

\choices
    {\(6\)}
    {\(2\)}
    {\(\frac{1}{2}\)}
    {\(\frac{2\sqrt{7}}{7}\)}
\end{problem}

\begin{answer}
B
\end{answer}

\begin{solution}
椭圆的长半轴为 \(m\)，焦半距满足 \(c=\sqrt{m^2-(m^2-1)}=1\). 由椭圆定义，点 \(P\) 到两焦点的距离之和为 \(2m\)，所以 \(2m=3+1=4\)，即 \(m=2\). 因而离心率为 \(e=\frac{c}{m}=\frac12\). 椭圆上点到右焦点的距离等于离心率乘以该点到右准线的距离，即 \(1=e\,d=\frac12d\)，所以 \(d=2\)，故选 B.
\end{solution}

\begin{problem}
设集合 \( S = \{x \mid |x - 2| > 3\} \)，\( T = \{x \mid a < x < a + 8\} \)，\( S \cup T = \R \)，则 \( a \) 的取值范围是（\(\quad\)）

\choices
    {\(-3 < a < -1\)}
    {\(-3 \leqslant a \leqslant -1\)}
    {\(a \leqslant -3\) 或 \(a \geqslant -1\)}
    {\(a < -3\) 或 \(a > -1\)}
\end{problem}

\begin{answer}
A
\end{answer}

\begin{solution}
由 \(\lvert x-2\rvert>3\)，得 \(S=(-\infty,-1)\cup(5,+\infty)\). 因此区间 \([-1,5]\) 中的每个点都必须属于 \(T=(a,a+8)\). 由于 \(T\) 为开区间，必须有 \(a<-1\) 且 \(a+8>5\)，即 \(-3<a<-1\).
反之，满足该条件时 \([-1,5]\subset T\)，从而 \(S\cup T=\R\). 故选 A.
\end{solution}

\begin{problem}
设函数 \( f(x) = \frac{1}{1 - \sqrt{x}} \)（\(0 \leqslant x < 1\)）的反函数为 \( f^{-1}(x) \)，则（\(\quad\)）

\choices
    {\(f^{-1}(x)\) 在其定义域上是增函数且最大值为\(1\)}
    {\(f^{-1}(x)\) 在其定义域上是减函数且最小值为\(0\)}
    {\(f^{-1}(x)\) 在其定义域上是减函数且最大值为\(1\)}
    {\(f^{-1}(x)\) 在其定义域上是增函数且最小值为\(0\)}
\end{problem}

\begin{answer}
D
\end{answer}

\begin{solution}
令 \(t=\sqrt{x}\)，则 \(t\in[0,1)\)，且 \(f(x)=\frac{1}{1-t}\).
当 \(t\) 增大时，\(1-t\) 减小，故 \(f(x)\) 在 \([0,1)\) 上递增，值域为 \([1,+\infty)\). 由反函数的性质，\(f^{-1}(x)\) 的定义域为 \([1,+\infty)\)，值域为 \([0,1)\)，且仍为增函数，最小值为 \(0\). 故选 D.
\end{solution}

\begin{problem}
已知函数 \( f(x) = \begin{cases}
-x + 1, & x < 0,\\
x - 1, & x \geqslant 0,
\end{cases} \) 则不等式 \( x + (x + 1)f(x + 1) \leqslant 1 \) 的解集是（\(\quad\)）

\choices
    {\(\{x\mid -1\leqslant x\leqslant \sqrt{2} -1\}\)}
    {\(\{x\mid x\leqslant 1\}\)}
    {\(\{x\mid x\leqslant \sqrt{2} -1\}\)}
    {\(\{x\mid -\sqrt{2} -1\leqslant x\leqslant \sqrt{2} -1\}\)}
\end{problem}

\begin{answer}
C
\end{answer}

\begin{solution}
当 \(x<-1\) 时，\(x+1<0\)，所以 \(f(x+1)=-x\). 原不等式化为
\(x-x(x+1)\leqslant1\iff -x^2\leqslant1\)，
对一切 \(x<-1\) 都成立.

当 \(x\geqslant-1\) 时，\(x+1\geqslant0\)，所以 \(f(x+1)=x\). 原不等式化为
\(x+x(x+1)\leqslant1\iff x^2+2x-1\leqslant0\)，解得 \(-1-\sqrt{2}\leqslant x\leqslant-1+\sqrt{2}\). 与 \(x\geqslant-1\) 联立，得 \(-1\leqslant x\leqslant\sqrt{2}-1\). 合并两种情况，解集为 \(x\leqslant\sqrt{2}-1\)，故选 C.
\end{solution}

\begin{problem}
已知函数 \( f(x) \) 是 \(\R\) 上的偶函数，且在区间 \([0, +\infty)\) 上是增函数. 令 \(a = f\left(\sin \frac{2\pi}{7}\right)\)，\(b = f\left(\cos \frac{5\pi}{7}\right)\)，\(c = f\left(\tan \frac{5\pi}{7}\right)\)，则（\(\quad\)）

\choices
    {\(b < a < c\)}
    {\(c < b < a\)}
    {\(b < c < a\)}
    {\(a < b < c\)}
\end{problem}

\begin{answer}
A
\end{answer}

\begin{solution}
因为 \(f\) 是偶函数，所以只需比较自变量的绝对值. 由 \(\left\lvert\cos\frac{5\pi}{7}\right\rvert=\cos\frac{2\pi}{7}\)、\(\left\lvert\tan\frac{5\pi}{7}\right\rvert=\tan\frac{2\pi}{7}\)，
且 \(\frac{2\pi}{7}>\frac\pi4\)，可得
\(\cos\frac{2\pi}{7}<\sin\frac{2\pi}{7}<\tan\frac{2\pi}{7}\).
由于 \(f\) 在 \([0,+\infty)\) 上递增，所以 \(b<a<c\)，故选 A.
\end{solution}

\begin{problem}
有 \(8\) 张卡片分别标有数字 \(1\)，\(2\)，\(3\)，\(4\)，\(5\)，\(6\)，\(7\)，\(8\)，从中取出 \(6\) 张卡片排成 \(3\) 行 \(2\) 列，要求 \(3\) 行中仅有中间行的两张卡片上的数字之和为 \(5\)，则不同的排法共有（\(\quad\)）

\choices
    {\(1344\) 种}
    {\(1248\) 种}
    {\(1056\) 种}
    {\(960\) 种}
\end{problem}

\begin{answer}
B
\end{answer}

\begin{solution}
中间行的两张卡片之和为 \(5\)，其数字只能是 \(\{1,4\}\) 或 \(\{2,3\}\)，共有 \(2\) 种选法，且在中间行的两个位置上有 \(2\) 种排列.

固定中间行后，余下的 \(6\) 张卡片中仍有且只有一对数字之和为 \(5\). 余下四个位置的总排法数为 \(\mathrm A_6^4=360\). 若上行或下行的两张卡片组成这对数字，则分别有 \(2\mathrm A_4^2=24\) 种排法；这两类不会同时发生. 因而固定中间行时，合乎“仅有中间行”的排法数为
\(360-24-24=312\). 总排法数为 \(2\times2\times312=1248\)，故选 B.
\end{solution}

\section{填空题}

\begin{problem}
\(\left(x - \frac{2}{\sqrt{x}}\right)^5\) 的二项展开式中，\(x^2\) 的系数是\fillinblank{}.（用数字作答）
\end{problem}

\begin{answer}
\(40\)
\end{answer}

\begin{solution}
在展开式中，若第二项取 \(k\) 次，则对应项为 \(\mathrm C_5^k x^{5-k}\left(-\frac{2}{\sqrt{x}}\right)^k=\mathrm C_5^k(-2)^k x^{5-\frac32k}\). 令幂指数满足 \(5-\frac32k=2\)，得 \(k=2\). 因此 \(x^2\) 的系数为 \(\mathrm C_5^2(-2)^2=10\times4=40\).
\end{solution}

\begin{problem}
一个正方体的各定点均在同一球的球面上，若该球的体积为 \(4\sqrt{3}\pi\)，则该正方体的表面积为\fillinblank{}.
\end{problem}

\begin{answer}
\(24\)
\end{answer}

\begin{solution}
设正方体的棱长为 \(s\)，外接球半径为 \(R\). 正方体的体对角线为 \(s\sqrt{3}\)，所以
\(R=\frac{s\sqrt{3}}2\). 由球体积为 \(4\sqrt{3}\pi\)，得 \(\frac43\pi R^3=4\sqrt{3}\pi\iff R^3=3\sqrt{3}\iff R=\sqrt{3}\). 于是 \(s=2\)，正方体的表面积为 \(6s^2=24\).
\end{solution}

\begin{problem}
已知圆 \(C\) 的圆心与抛物线 \(y^{2} = 4x\) 的焦点关于直线 \(y = x\) 对称. 直线 \(4x - 3y - 2 = 0\) 与圆 \(C\) 相交于 \(A\)，\(B\) 两点，且 \(|AB| = 6\)，则圆 \(C\) 的方程为\fillinblank{}.
\end{problem}

\begin{answer}
\(x^{2} + (y - 1)^{2} = 10\)
\end{answer}

\begin{solution}
抛物线 \(y^2=4x\) 的焦点为 \((1,0)\)，关于直线 \(y=x\) 对称后，圆 \(C\) 的圆心为 \((0,1)\). 圆心到直线 \(4x-3y-2=0\) 的距离为 \(d=\frac{\lvert4\times0-3\times1-2\rvert}{\sqrt{4^2+(-3)^2}}=1\). 设圆半径为 \(r\)，则弦长公式给出 \(6=2\sqrt{r^2-d^2}=2\sqrt{r^2-1}\)，从而 \(r^2=10\). 所以圆 \(C\) 的方程为 \(x^2+(y-1)^2=10\).
\end{solution}

\begin{problem}
如图，在平行四边形 \(ABCD\) 中，\(\overrightarrow{AC} = (1, 2)\)，\(\overrightarrow{BD} = (-3, 2)\)，则 \(\overrightarrow{AD} \cdot \overrightarrow{AC} = \)\fillinblank{}.

\bitmapfigure[width=0.42\linewidth]{img/2008/tianjin_science/q14_fig1.png}
\end{problem}

\begin{answer}
\(3\)
\end{answer}

\begin{solution}
设 \(\overrightarrow{AB}=\bs{u}\)，\(\overrightarrow{AD}=\bs{v}\). 由平行四边形的对角线和对角线差，得 \(\bs{u}+\bs{v}=\overrightarrow{AC}=(1,2)\)，\(\bs{v}-\bs{u}=\overrightarrow{BD}=(-3,2)\). 两式相加得 \(2\bs{v}=(-2,4)\)，所以 \(\overrightarrow{AD}=\bs{v}=(-1,2)\). 因而 \(\overrightarrow{AD}\cdot\overrightarrow{AC}=(-1,2)\cdot(1,2)=-1+4=3\).
\end{solution}

\begin{problem}
已知数列 \( \{a_{n}\} \) 中， \( a_{1}=1 \)， \( a_{n+1}-a_{n}=\frac{1}{3^{n+1}} \)（\(n\in \N^{*}\)），则 \( \displaystyle\lim_{n\to\infty}a_{n}= \)\fillinblank{}.
\end{problem}

\begin{answer}
\(\frac{7}{6}\)
\end{answer}

\begin{solution}
由递推关系，从 \(a_1\) 累加得到 \(a_n=1+\sum_{k=1}^{n-1}\frac{1}{3^{k+1}}=1+\sum_{k=2}^{n}\frac{1}{3^k}\). 因此 \(\lim_{n\to\infty}a_n=1+\sum_{k=2}^{\infty}\frac{1}{3^k}=1+\frac{1/9}{1-1/3}=\frac76\).
\end{solution}

\begin{problem}
设 \(a > 1\)，若仅有一个常数 \(c\) 使得对于任意的 \(x \in [a, 2a]\)，都有 \(y \in [a, a^2]\) 满足方程 \(\log_a x + \log_a y = c\)，这时，\(a\) 的取值的集合为\fillinblank{}.
\end{problem}

\begin{answer}
\(\{2\}\)
\end{answer}

\begin{solution}
令 \(r=\log_a2\). 因为 \(a>1\)，所以 \(r>0\). 对固定的 \(x\in[a,2a]\)，存在 \(y\in[a,a^2]\) 使等式成立，当且仅当
\(c=\log_a x+\log_a y\in[\log_a x+1,\log_a x+2]\).
要使这个条件对所有 \(x\in[a,2a]\) 成立，常数 \(c\) 必须属于所有这些区间的交集. 其中左端点的最大值在 \(x=2a\) 处取得，为 \(2+r\)，右端点的最小值在 \(x=a\) 处取得，为 \(3\)，所以所有可能的常数构成区间 \([2+r,3]\). 当 \(2+r<3\) 时有无穷多个常数，当 \(2+r>3\) 时不存在这样的常数；题设要求仅有一个常数，故必须有 \(2+r=3\)，即 \(\log_a2=1\)，从而 \(a=2\). 因此 \(a\) 的取值集合为 \(\{2\}\).
\end{solution}

\section{解答题}

\begin{problem}
已知 \(\cos \left(x - \frac{\pi}{4}\right) = \frac{\sqrt{2}}{10}\)，\(x \in \left(\frac{\pi}{2}, \frac{3\pi}{4}\right)\).
\begin{enumerate}
\item 求 \(\sin x\) 的值；
\item 求 \( \sin\left(2x+\frac{\pi}{3}\right) \) 的值.
\end{enumerate}
\end{problem}

\begin{solution}
\begin{enumerate}
\item
令 \(u=x-\frac{\pi}{4}\)，则 \(u\in\left(\frac{\pi}{4},\frac{\pi}{2}\right)\)，从而 \(\sin u>0\). 由题设，
\(\sin u=\sqrt{1-\cos^2u}=\sqrt{1-\frac{2}{100}}=\frac{7\sqrt{2}}{10}\).

因此 \(\sin x=\sin\left(u+\frac{\pi}{4}\right)=\sin u\cos\frac{\pi}{4}+\cos u\sin\frac{\pi}{4}=\frac{7}{10}+\frac{1}{10}=\frac45\).

因此 \(\sin x=\frac45\).

\item
又因为 \(\cos x=\cos\left(u+\frac{\pi}{4}\right)=\cos u\cos\frac{\pi}{4}-\sin u\sin\frac{\pi}{4}=\frac{1}{10}-\frac{7}{10}=-\frac35\)，所以 \(\sin 2x=2\sin x\cos x=-\frac{24}{25}\)，\(\cos 2x=\cos^2x-\sin^2x=-\frac{7}{25}\). 从而 \(\sin\left(2x+\frac{\pi}{3}\right)=\sin 2x\cos\frac{\pi}{3}+\cos 2x\sin\frac{\pi}{3}=-\frac{12}{25}-\frac{7\sqrt{3}}{50}=-\frac{24+7\sqrt{3}}{50}\).
\end{enumerate}
\end{solution}

\begin{problem}
甲，乙两个篮球运动员互不影响地在同一位置投球，命中率分别为 \(\frac{1}{2}\) 与 \(p\)，且乙投球 \(2\) 次均未命中的概率为 \(\frac{1}{16}\).
\begin{enumerate}
\item 求乙投球的命中率 \(p\).
\item 若甲投球 \(1\) 次，乙投球 \(2\) 次，两人共命中的次数记为 \(\xi\)，求 \(\xi\) 的分布列和数学期望.
\end{enumerate}
\end{problem}

\begin{solution}
\begin{enumerate}
\item
设乙两次投球命中的事件分别为 \(B_1\)、\(B_2\). 由两次投球相互独立，\(P(\overline{B_1}\cap\overline{B_2})=(1-p)^2=\frac1{16}\).
因为 \(0\leqslant p\leqslant1\)，所以 \(1-p=\frac14\)，故 \(p=\frac34\).

\item
记甲投球命中为事件 \(A\)，乙两次投球命中的次数分别记为 \(Y_1\)，\(Y_2\). 则甲命中概率为 \(\frac12\)，乙每次命中概率为 \(\frac34\)，且三次投球相互独立，\(\xi\) 的可能取值为 \(0\)，\(1\)，\(2\)，\(3\). 因此
\(P(\xi=0)=\frac12\cdot\frac14\cdot\frac14=\frac1{32}\). 当 \(\xi=1\) 时，有“甲未命中且乙恰命中一次”或“甲命中且乙两次均未命中”两种情形，所以 \(P(\xi=1)=\frac12\cdot2\cdot\frac34\cdot\frac14+\frac12\cdot\frac14\cdot\frac14=\frac7{32}\). 当 \(\xi=2\) 时，有“甲未命中且乙两次均命中”或“甲命中且乙恰命中一次”两种情形，所以 \(P(\xi=2)=\frac12\cdot\frac34\cdot\frac34+\frac12\cdot2\cdot\frac34\cdot\frac14=\frac{15}{32}\). 最后 \(P(\xi=3)=\frac12\cdot\frac34\cdot\frac34=\frac9{32}\).

分布列为
\begin{center}
\begin{tabular}{c|cccc}
\hline
\(\xi\) & \(0\) & \(1\) & \(2\) & \(3\) \\
\hline
\(P\) & \(\frac1{32}\) & \(\frac7{32}\) & \(\frac{15}{32}\) & \(\frac9{32}\) \\
\hline
\end{tabular}
\end{center}

所以
\(E(\xi)=0\cdot\frac1{32}+1\cdot\frac7{32}+2\cdot\frac{15}{32}+3\cdot\frac9{32}=2\).
\end{enumerate}
\end{solution}

\begin{problem}
如图，在四棱锥 \(P - ABCD\) 中，底面 \(ABCD\) 是矩形. 已知 \(AB = 3\)，\(AD = 2\)，\(PA = 2\)，\(PD = 2\sqrt{2}\)，\(\angle PAB = 60^{\circ}\).

\begin{enumerate}
\item 证明 \(AD \perp\) 平面 \(PAB\)；
\item 求异面直线 \(PC\) 与 \(AD\) 所成角的大小；
\item 求二面角 \(P - BD - A\) 的大小.
\end{enumerate}

\bitmapfigure[width=0.34\linewidth]{img/2008/tianjin_science/q19_fig1.png}
\end{problem}

\begin{solution}
\begin{enumerate}
\item
在 \(\triangle PAD\) 中，
\(PA^2+AD^2=2^2+2^2=8=(2\sqrt{2})^2=PD^2\)，
所以 \(PA\perp AD\). 又因为 \(ABCD\) 是矩形，\(AB\perp AD\). 直线 \(PA\)，\(AB\) 是平面 \(PAB\) 内相交的两条直线，故 \(AD\perp\) 平面 \(PAB\).

\item
因为 \(AD\parallel BC\)，所以异面直线 \(PC\) 与 \(AD\) 所成的角等于直线 \(PC\) 与 \(BC\) 所成的锐角，记为 \(\theta\). 在 \(\triangle PAB\) 中，由余弦定理，
\(PB^2=PA^2+AB^2-2PA\cdot AB\cos\angle PAB=4+9-12\cdot\frac12=7\)，
即 \(PB=\sqrt{7}\). 由 \(AD\perp\) 平面 \(PAB\)，得 \(AD\perp PB\)，再由 \(BC\parallel AD\)，得 \(BC\perp PB\). 因而 \(\triangle PBC\) 在 \(B\) 处为直角三角形，且 \(BC=AD=2\). 所以
\(\tan\theta=\frac{PB}{BC}=\frac{\sqrt{7}}{2}\)，故 \(\theta=\arctan\frac{\sqrt{7}}{2}\).

\item
过点 \(P\) 作 \(PH\perp AB\)，垂足为 \(H\)，再过点 \(H\) 作 \(HE\perp BD\)，垂足为 \(E\). 由 \(AD\perp\) 平面 \(PAB\)，可知 \(PH\perp AD\)，又 \(PH\perp AB\)，所以 \(PH\perp\) 平面 \(ABCD\). 因而 \(HE\) 是 \(PE\) 在平面 \(ABCD\) 内的射影，且由三垂线定理，\(PE\perp BD\). 所以 \(\angle PEH\) 是二面角 \(P-BD-A\) 的平面角.

由题设，
\(PH=PA\sin60^{\circ}=\sqrt{3}\)，\(AH=PA\cos60^{\circ}=1\)，从而 \(BH=AB-AH=2\). 在直角三角形 \(ABD\) 中，\(BD=\sqrt{AB^2+AD^2}=\sqrt{13}\). 因为 \(HE\) 是点 \(H\) 到直线 \(BD\) 的距离，利用三角形面积可得 \(HE=\frac{BH\cdot AD}{BD}=\frac4{\sqrt{13}}\). 在直角三角形 \(PHE\) 中，\(\tan\angle PEH=\frac{PH}{HE}=\frac{\sqrt{3}}{4/\sqrt{13}}=\frac{\sqrt{39}}4\)，
故二面角 \(P-BD-A\) 的大小为 \(\arctan\frac{\sqrt{39}}4\).
\end{enumerate}
\end{solution}

\begin{problem}
已知函数 \( f(x) = x + \frac{a}{x} + b \)（\(x \neq 0\)），其中 \(a\)，\(b \in \R\).
\begin{enumerate}
\item 若曲线 \( y = f(x) \) 在点 \( P(2, f(2)) \) 处的切线方程为 \( y = 3x + 1 \)，求函数 \( f(x) \) 的解析式；
\item 讨论函数 \( f(x) \) 的单调性；
\item 若对于任意的 \(a \in \left[\frac{1}{2}, 2\right]\)，不等式 \(f(x) \leqslant 10\) 在 \(\left[\frac{1}{4}, 1\right]\) 上恒成立，求 \(b\) 的取值范围.
\end{enumerate}
\end{problem}

\begin{solution}
\begin{enumerate}
\item
函数的定义域为 \(x\neq0\)，且
\(f'(x)=1-\frac{a}{x^2}=\frac{x^2-a}{x^2}\).

在点 \(P(2,f(2))\) 处的切线斜率为 \(3\)，所以
\(f'(2)=1-\frac{a}{4}=3\)，
解得 \(a=-8\). 又因为点 \(P\) 在直线 \(y=3x+1\) 上，故 \(f(2)=7\). 于是
\(f(2)=2-\frac82+b=b-2=7\)，解得 \(b=9\). 因此函数的解析式为 \(f(x)=x-\frac8x+9\)，\((x\neq0)\).

\item
当 \(a\leqslant0\) 时，对任意 \(x\neq0\)，有 \(x^2-a>0\)，所以 \(f'(x)>0\). 因此 \(f(x)\) 在 \(\left(-\infty,0\right)\) 和 \((0,+\infty)\) 上都是增函数.

当 \(a>0\) 时，令 \(f'(x)=0\)，得 \(x=\pm\sqrt{a}\). 由 \(f'(x)\) 的符号可知，\(f(x)\) 在 \(\left(-\infty,-\sqrt{a}\right)\) 和 \((\sqrt{a},+\infty)\) 上是增函数，在 \((-\sqrt{a},0)\) 和 \((0,\sqrt{a})\) 上是减函数.

\item
第三问中，参数 \(a\in\left[\frac12,2\right]\)，且 \(x\in\left[\frac14,1\right]\). 对固定的正数 \(x\)，\(x+\frac{a}{x}+b\) 随 \(a\) 增大而增大，所以只需取 \(a=2\). 此时
\(g(x)=x+\frac2x+b\)，且 \(g'(x)=1-\frac2{x^2}<0\)，\(\left(x\in\left[\frac14,1\right]\right)\). 故 \(g(x)\) 在该区间上递减，最大值在 \(x=\frac14\) 处取得. 由题意，\(\frac14+\frac{2}{1/4}+b=\frac{33}{4}+b\leqslant10\)，
所以 \(b\leqslant\frac74\). 因而 \(b\) 的取值范围为 \(\left(-\infty,\frac74\right]\).
\end{enumerate}
\end{solution}

\begin{problem}
已知中心在原点的双曲线 \(C\) 的一个焦点是 \(F_{1}(-3,0)\)，一条渐近线的方程是 \(\sqrt{5} x - 2y = 0\).
\begin{enumerate}
\item 求双曲线 \(C\) 的方程；
\item 若以 \(k\)（\(k \neq 0\)）为斜率的直线 \(l\) 与双曲线 \(C\) 相交于两个不同的点 \(M\)，\(N\)，线段 \(MN\) 的垂直平分线与两坐标轴围成的三角形的面积为 \(\frac{81}{2}\)，求 \(k\) 的取值范围.
\end{enumerate}
\end{problem}

\begin{solution}
\begin{enumerate}
\item
设双曲线的标准方程为 \(\frac{x^2}{A^2}-\frac{y^2}{B^2}=1\)，其中 \(A>0\)，\(B>0\). 其焦半距满足 \(A^2+B^2=3^2=9\). 又因为渐近线的斜率为 \(\frac{B}{A}\)，由题设渐近线 \(\sqrt{5}x-2y=0\) 得
\(\frac{B}{A}=\frac{\sqrt{5}}2\)，\(B^2=\frac54A^2\). 所以 \(A^2=4\)，\(B^2=5\)，双曲线 \(C\) 的方程为 \(\frac{x^2}{4}-\frac{y^2}{5}=1\).

\item
设直线 \(l\) 的方程为 \(y=kx+m\)，其中 \(k\neq0\)，并记 \(d=5-4k^2\). 将直线方程代入双曲线方程，得
\(d x^2-8kmx-4m^2-20=0\).
因为直线与双曲线有两个不同的交点，所以 \(d\neq0\)，且该方程的判别式为正. 由根与系数的关系，线段 \(MN\) 的中点 \(Q(x_0,y_0)\) 满足
\(x_0=\frac{4km}{d}\)，\(y_0=kx_0+m=\frac{5m}{d}\).

直线 \(MN\) 的斜率为 \(k\)，所以其垂直平分线经过 \(Q\)，斜率为 \(-\frac1k\)，方程为 \(x+ky=x_0+ky_0=\frac{9km}{d}\).
它在 \(x\) 轴、\(y\) 轴上的截距分别为 \(\frac{9km}{d}\)、\(\frac{9m}{d}\). 由三角形面积为 \(\frac{81}{2}\)，得
\(\frac12\left|\frac{9km}{d}\cdot\frac{9m}{d}\right|=\frac{81}{2}\)，即 \(m^2=\frac{d^2}{|k|}\).

原二次方程的判别式为
\[
\Delta=(-8km)^2-4d(-4m^2-20)=80\left(m^2+5-4k^2\right)=80\left(m^2+d\right)
\]
因此还需满足 \(m^2+d>0\). 令 \(t=|k|>0\)，则 \(d=5-4t^2\)，代入上式得到
\(m^2+d=\frac{d^2}{t}+d=\frac{d(d+t)}{t}>0\).

当 \(0<t<\frac{\sqrt{5}}2\) 时，\(d>0\)，上式自动成立；当 \(t>\frac{\sqrt{5}}2\) 时，\(d<0\)，需有 \(d+t<0\)，即
\(5-4t^2+t<0\iff 4t^2-t-5>0\iff t>\frac54\).
故
\(|k|\in\left(0,\frac{\sqrt{5}}2\right)\cup\left(\frac54,+\infty\right)\).
所以 \(k\) 的取值范围为
\[
\left(-\infty,-\frac54\right)\cup\left(-\frac{\sqrt{5}}2,0\right)\cup\left(0,\frac{\sqrt{5}}2\right)\cup\left(\frac54,+\infty\right)
\]
\end{enumerate}
\end{solution}

\begin{problem}
在数列 \(\{a_{n}\}\) 与 \(\{b_{n}\}\) 中，\(a_{1} = 1\)，\(b_{1} = 4\)，数列 \(\{a_{n}\}\) 的前 \(n\) 项和 \(S_{n}\) 满足 \(nS_{n + 1} - (n + 3)S_{n} = 0\)，\(2a_{n + 1}\) 为 \(b_{n}\) 与 \(b_{n + 1}\) 的等比中项，\(n \in \N^{*}\).
\begin{enumerate}
\item 求 \(a_{2}\)，\(b_{2}\) 的值；
\item 求数列 \(\{a_{n}\}\) 与 \(\{b_{n}\}\) 中的通项公式；
\item 设 \(T_{n} = (-1)^{a_{1}}b_{1} + (-1)^{a_{2}}b_{2} + \dots + (-1)^{a_{n}}b_{n}\)，\(n \in \N^{*}\)；证明 \(|T_{n}| < 2n^{2}\)，\(n \geqslant 3\).
\end{enumerate}
\end{problem}

\begin{solution}
\begin{enumerate}
\item
取 \(n=1\) 代入前 \(n\) 项和的关系，得
\(S_2-4S_1=0\).
因为 \(S_1=a_1=1\)，所以 \(S_2=4\)，从而 \(a_2=S_2-S_1=3\). 又由等比中项的定义，
\((2a_2)^2=b_1b_2\)，
故 \(36=4b_2\)，即 \(b_2=9\).

\item
当 \(n\geqslant2\) 时，题设关系分别写成 \(nS_{n+1}=(n+3)S_n\)，\((n-1)S_n=(n+2)S_{n-1}\).
由第一式得 \(na_{n+1}=3S_n\)，由第二式得 \((n-1)a_n=3S_{n-1}\). 因而
\(na_{n+1}=3S_n=3S_{n-1}+3a_n=(n+2)a_n\)，即 \(a_{n+1}=\frac{n+2}{n}a_n\). 结合 \(a_1=1\)，\(a_2=3\)，逐项相乘可得 \(a_n=3\prod_{j=2}^{n-1}\frac{j+2}{j}=3\cdot\frac{4\cdot5\cdots(n+1)}{2\cdot3\cdots(n-1)}=\frac{n(n+1)}2\). 该式在 \(n=1\) 时也成立，所以 \(a_n=\frac{n(n+1)}2\).

由等比中项条件，\(b_nb_{n+1}=(2a_{n+1})^2=\left((n+1)(n+2)\right)^2\). 令 \(x_n=\frac{b_n}{(n+1)^2}\)，则 \(x_1=1\)，且 \(x_nx_{n+1}=\frac{b_nb_{n+1}}{(n+1)^2(n+2)^2}=1\). 由 \(x_1=1\) 递推得到所有 \(x_n=1\)，故 \(b_n=(n+1)^2\).

\item
最后证明不等式. 因为 \(a_n=\frac{n(n+1)}2\)，
所以 \(a_n\) 的奇偶性按 \(n=1\)，\(2\)，\(3\)，\(4\) 为奇、奇、偶、偶循环. 因而 \((-1)^{a_n}\) 按四项一组为负、负、正、正. 又 \(b_n=(n+1)^2\)，于是按 \(n\) 的模 \(4\) 分类求和，若 \(k\in\N^{*}\)，则
\[
T_{4k}=\sum_{j=0}^{k-1}\left[-(4j+2)^2-(4j+3)^2+(4j+4)^2+(4j+5)^2\right]=16k^2+12k
\]
\[
T_{4k+1}=T_{4k}-(4k+2)^2=-4k-4
\]
\[
T_{4k+2}=T_{4k+1}-(4k+3)^2=-16k^2-28k-13
\]
\[
T_{4k+3}=T_{4k+2}+(4k+4)^2=4k+3
\]

当 \(n=4k\) 时，\(k\geqslant1\)，且
\[
|T_n|=16k^2+12k<32k^2=2n^2
\]
当 \(n=4k+1\) 时，\(k\geqslant1\)，且
\[
|T_n|=4k+4<2(4k+1)^2=2n^2
\]
当 \(n=4k+2\) 时，\(k\geqslant1\)，且
\[
|T_n|=16k^2+28k+13<32k^2+32k+8=2(4k+2)^2=2n^2
\]
当 \(n=4k+3\) 时，\(k\geqslant0\)，且
\[
|T_n|=4k+3<2(4k+3)^2=2n^2
\]
综上，当 \(n\geqslant3\) 时，恒有 \(|T_n|<2n^2\).
\end{enumerate}
\end{solution}
